\section{Theoretische Grundlagen}

\subsection{Software-Testing und Testfälle}
Der Bereich des Software-Testings beschäftigt sich damit, Systeme und Programme 
auf ihre Korrektheit, Funktionalität, Vollständigkeit, Sicherheit und generelle Qualität 
zu überprüfen und gilt dabei als wesentliche Voraussetzung für die 
Qualität und Zuverlässigkeit von Softwareprodukten \cite{ehmer_comparative_2012, wang_software_2024}.
Allgemeiner formuliert besteht das Ziel von Software-Testing darin, 
Abweichungen zwischen dem spezifizierten und dem tatsächlichen Verhalten eines Systems zu 
identifizieren. Programme werden also unter festgelegten Bedingungen ausgeführt, wodurch 
geprüft werden kann, ob die vorab definierten Anforderungen erfüllt werden oder ob Fehler 
auftreten.

Konzeptionell lassen sich die verschiedenen Arten und Techniken zum Testen von Software in die Kategorien White-Box-Testing 
und Black-Box-Testing einteilen \cite{nidhra_black_2012}. Für das White-Box-Testing benötigt 
der Tester und Testfallersteller Zugang zum Quellcode des Programms. Die Tests überprüfen primär die 
interne Logik und Struktur des Codes, indem Funktionen mit unterschiedlichen Parametern ausgeführt werden, um 
möglichst viele logische Pfade im Programmablauf auf ihre Funktionalität untersuchen zu können \cite{verma_comparative_2017}.

Das Black-Box-Testing, auch Functional Testing genannt, behandelt den Quellcode des zu testenden Programms 
hingegen wie eine undurchsichtige Box, deren Inhalt nicht bekannt ist und auch nicht bekannt sein soll. 
Ohne Wissen über die internen Abläufe werden Testfälle und erwartete Ergebnisse nur aus Spezifikationen 
oder Beschreibungen über das gewünschte Verhalten eines Programms abgeleitet \cite{verma_comparative_2017}. Auf Basis dieser 
Testfälle wird anschließend geprüft, ob das Programm die definierte Funktionalität erfüllt.

Gerade beim Black-Box-Testing hängt die Qualität eines Testprozesses maßgeblich von der Auswahl geeigneter Testfälle ab. 
Nach der Definition des International Software Testing Qualifications Board besteht ein Testfall primär aus Vorbedingungen, 
Eingaben, gegebenenfalls auszuführenden Aktionen, erwarteten Ergebnissen und Nachbedingungen \cite{istqb_testfall}. 
Für die in dieser Projektarbeit betrachteten algorithmischen Programmieraufgaben lässt sich die allgemeine Definition 
vereinfachen. Jede Aufgabe definiert festgelegte Parameter und Rückgabewerte. Somit besteht ein Testfall immer aus einer 
konkreten Eingabe und dem dazugehörigen erwarteten Ergebnis. Das erwartete Ergebnis beschreibt das Verhalten, 
das gemäß Spezifikation bei den jeweiligen Testfällen eintreten soll \cite{istqb_erwartetes_ergebnis}.


\subsection{KI-gestützte Testfallgenerierung}
Diese Projektarbeit baut wesentlich auf dem Konzept des Black-Box-Testings auf. Vor dem Hintergrund, 
dass Large Language Models zunehmend zur Generierung von Quellcode eingesetzt werden, liegt es nahe, 
sie zugleich auch Tests und Testfälle erstellen zu lassen.
So werden LLMs inzwischen auch für die Vorbereitung von
Testfällen, die Erzeugung von Unit-Tests und Testorakeln sowie die Generierung von Eingaben für
Systemtests verwendet \cite{wang_software_2024}.
Dabei entsteht jedoch ein methodisches Risiko, wenn die Implementierung und Tests durch dasselbe 
Modell erzeugt werden. Hat das Modell die spezifizierte Funktionalität bereits im Code falsch 
umgesetzt, kann nicht ausgeschlossen werden, dass es anschließend Testfälle mit erwarteten 
Ergebnissen generiert, die auf derselben Fehlinterpretation beruhen. In diesem Fall würden die Tests 
nicht das ursprünglich geforderte Verhalten prüfen, sondern lediglich bestätigen, dass der erzeugte 
Code konsistent mit der eigenen, möglicherweise fehlerhaften Auslegung der Aufgabe ist \cite{caridad_modelinversion_2026}. 
Vor diesem Hintergrund erscheint es besonders relevant zu untersuchen, wie zuverlässig ein 
LLM Testfälle ausschließlich aus der natürlichsprachlichen Spezifikation ableiten 
kann, ohne Zugriff auf die konkrete Implementierung zu erhalten.

Gerade aus diesem Grund ist für diese Untersuchung die Zweiteiligkeit eines Testfalls, bestehend aus Eingabe 
und erwartetem Ergebnis besonders relevant. Beide Angaben werden von dem verwendeten LLM 
auf Basis einer natürlichsprachlichen Programmieraufgabe (Spezifikation des Sollverahltens) erzeugt.
Die empirische Forschung zeigt jedoch, dass die Korrektheit des vom 
LLM generierten erwarteten Ergenisses je nach Aufgabe und Untersuchungsaufbau deutlich variieren kann
\cite{wang_software_2024}. Die Korrektheit hängt davon ab, wie gut das Modell die Programmieraufgabe 
verstanden hat und wie gut es damit in der Lage war, das Soll-Ergebnis der eigenen Testfälle vorherzusagen.
Ein KI-generierter Testfall wird in dieser Arbeit demnach erst als valide eingestuft, wenn die 
Eingabe die jeweiligen Einschränkungen innerhalb der Programmieraufgabe erfüllt und das vom LLM 
erwartete Ergebnis mit dem Ergebnis der korrekten Referenzimplementierung übereinstimmt. 
Erst dann kann untersucht werden, ob die Gruppe der validierten Testfälle auch dazu geeignet ist, 
bewusst fehlerhafte Implementierungen zu identifizieren.

\subsection{Large Language Models: Funktionsweise und Auswahl}
Um einordnen zu können, weshalb ein Large Language Model sowohl geeignete Testfälle als auch fehlerhafte erwartete 
Ergebnisse erzeugen kann, ist zunächst ein grundlegendes Verständnis seiner Funktionsweise erforderlich. 

Moderne Large Language Models basieren überwiegend auf der von Vaswani et al. vorgestellten 
Transformer-Architektur. Hierbei ermöglichen sogenannte Attention-Mechanismen unterschiedliche Stellen 
in einer sprachlichen Sequenz miteinander in Beziehung zu setzen \cite{vaswani_attention_nodate}. 
Auf dieses Projekt bezogen, können somit Zusammenhänge zwischen den Bedingungen einer natürlichsprachlichen 
Programmieraufgaben, den darin definierten Eingabegrenzen und dem geforderten 
Rückgabewert modelliert und als Kontext während der Generierung der Antwort verwendet werden.

Die Ausgabe eines Sprachmodells wird schrittweise in sogenannten Tokens erzeugt. Tokens stellen abhängig von der verwendeten 
Tokenisierung vollständige Wörter oder Wortteile dar. Auf Basis der Eingabe und der bereits generierten Tokens berechnet 
das Modell eine Wahrscheinlichkeitsverteilung für das jeweils nächste Token \cite{vaswani_attention_nodate}. 
Die Ausgabe eines generativen Sprachmodells stellt folglich keine bewiesen richtigen Fakten dar. 
Das Modell generiert eine sprachlich und strukturell wahrscheinliche Fortsetzung auf Grundlage der bisher verarbeiteten 
Tokens und des gesammelten Kontexts. Bezogen auf die Generierung von Testfällen erklärt 
dieses Verhalten, warum das Modell die Bedingungen und Vorgaben innerhalb einer Programmieraufgabe 
zwar in Verbindung setzen und als Kontext in die Ausgabe mit einfließen lassen kann, es jedoch nicht 
garantiert ist, dass die Ausgabe alle Vorgaben befolgt und fachlich korrekt ist \cite{huang_survey_2025}.  

Für die Umsetzung des zu dieser Projektarbeit zugehörigen Assistenzsystems wird das von OpenAI bereitgestellte 
Modell GPT-4.1 mini verwendet. Die Wahl dieses Modells 
beruht auf einem Kompromiss zwischen Leistungsfähigkeit, Antwortgeschwindigkeit und Nutzungskosten. 
Die Auswahl eines Mini-Modells ist bei der Auswertung und Interpretation der Ergebnisse zu beachten. 
Die Untersuchung zielt nicht darauf ab, die generell bestmögliche Testqualität eines aktuellen Sprachmodells zu erzielen. 
Vielmehr wird untersucht, welche Qualität mit der konkreten Kombination aus dem Modell GPT-4.1 mini und dem verwendeten System-Prompt 
bei den untersuchten Programmieraufgaben erreicht wird.

\newpage

